\documentclass[11pt,a4paper]{article}
\usepackage[margin=1in]{geometry}
\usepackage{times}
\usepackage{latexsym}
\usepackage{amsmath,amssymb,amsfonts}
\usepackage{booktabs}
\usepackage{multirow}
\usepackage{graphicx}
\usepackage{xcolor}
\usepackage{bm}
\usepackage{microtype}
\usepackage{natbib}
\usepackage{tabularx}
\usepackage{array}
\usepackage{hyperref}

\newcommand{\contamcal}{\textsc{ContamCal}}
\newcommand{\pp}{\,\text{pp}}

\title{\contamcal{}: Calibrating Benchmark Scores for Training Data Contamination\\via Contamination-Augmented Item Response Theory}

\author{Anonymous}

\begin{document}
\maketitle

%======================================================================
\begin{abstract}
%======================================================================
Training data contamination inflates LLM benchmark scores, yet existing detection methods flag contamination without answering: \emph{what would this model score if it had never seen the test data?}
We introduce \contamcal{}, a two-stage framework that bridges contamination detection and score calibration.
Stage~1 computes item-level contamination exposure via four complementary detection signals---embedding similarity, Min-K++, ConStat, and LLM self-critique---aggregated through Reciprocal Rank Fusion (RRF).
Stage~2 fits a contamination-augmented 4PL Item Response Theory model with per-benchmark contamination susceptibility parameters~$\gamma_k$ and a measurement error model that propagates signal uncertainty into Bayesian posteriors.
On controlled contamination injection experiments across Qwen2.5 (0.5B--7B) and Llama-3.2-1B with three benchmarks (MMLU, GSM8K, ARC-Challenge), \contamcal{} recovers true benchmark performance within 1.8\pp{} MAE (bootstrap 95\% CI: [1.2, 2.4]) with 94\% credible interval coverage, while naive scores overestimate by up to 58\pp{} under full contamination (3B GSM8K).
We systematically characterize signal effectiveness as method-specific and scale-dependent: embedding similarity is robust across scales (AUC$>$0.95), whereas Min-K++ and self-critique fail at 0.5B (AUC$\approx$0.47) but recover at 1.5B+ (AUC$>$0.70).
Fisher Information analysis establishes AUC$\geq$0.55 as the structural identifiability boundary and AUC$\geq$0.70 as the practical boundary for reliable calibration.
\end{abstract}

%======================================================================
\section{Introduction}
\label{sec:intro}
%======================================================================

Benchmark contamination---where test data leaks into training corpora---is among the most serious threats to the validity of LLM evaluation \citep{sainz2023contamination,balloccu2024leak,li2024task}.
Even rephrased test items can inflate accuracy by 8--15 points while evading $n$-gram decontamination \citep{yang2023rethinking,dekoninck2024evading}.
\citet{xu2024benchmark} systematically tested 20 mitigation strategies and found none reliably eliminates contamination effects.

This motivates a fundamental reframing: rather than trying to eliminate contamination from benchmarks (which does not work), we propose to \emph{statistically control for it}---measuring its effect and subtracting it out with proper uncertainty quantification.
This mirrors the shift in causal inference from eliminating confounders to statistically adjusting for them \citep{rosenbaum1983central}.

Existing approaches address only part of the problem.
Detection methods---Min-K\%++ \citep{shi2024detecting}, ConStat \citep{dekoninck2024constat}, and LLM self-critique \citep{yang2023rethinking}---flag contamination but do not calibrate scores.
IRT-based evaluation \citep{lalor2024benchmarking,polo2024tinybenchmarks} models item difficulty psychometrically but ignores contamination.
DCR \citep{dong2024generalization} provides adjusted accuracy but lacks uncertainty bounds.
HiBayES \citep{zhuang2024hibayes} provides Bayesian model evaluation but does not account for contamination.

We introduce \contamcal{}, which bridges detection and calibration through a two-stage pipeline grounded in the educational measurement tradition of item compromise IRT \citep{eckerly2017detecting,sinharay2017use} and explanatory IRT \citep{deboeck2004explanatory}.
Our framework makes three non-trivial adaptations for LLM evaluation: (a)~independent examinees become a single model with joint item distributions, altering local independence semantics; (b)~binary leakage becomes continuous exposure scores requiring measurement error propagation; and (c)~stable ability becomes regime-conditional ability requiring scale-aware signal selection.

Our contributions are:
\begin{enumerate}
\item The \contamcal{} framework: a two-stage pipeline (RRF provenance mapping $\to$ contamination-augmented 4PL IRT with per-benchmark $\gamma_k$ and measurement error model) producing calibrated scores with credible intervals (\S\ref{sec:method}).
\item Systematic signal effectiveness characterization: a signal $\times$ scale $\times$ dosage AUC mapping revealing method-specific, scale-dependent reliability, including a cautionary finding on Min-K++ directional reversal (\S\ref{sec:signal_analysis}).
\item Dual-track validation: controlled SFT injection across 0.5B--7B within Qwen2.5 plus Llama-3.2-1B (internal validity), and ecological observational studies on Instruct and DeepSeek-R1 models up to 72B (external validity) (\S\ref{sec:experiments}).
\end{enumerate}

%======================================================================
\section{Related Work}
\label{sec:related}
%======================================================================

\paragraph{Contamination Detection.}
Several methods detect training data memorization in LLMs.
Min-K\%++ \citep{shi2024detecting} uses token-level probability outliers; ConStat \citep{dekoninck2024constat} exploits performance gaps between original and paraphrased items; KDS \citep{zhang2024kds} and CDD \citep{li2024cdd} provide detection datasets.
\citet{yang2023rethinking} propose LLM self-critique for contamination detection.
These methods flag items but do not produce calibrated scores.

\paragraph{IRT for LLM Evaluation.}
Item Response Theory has been applied to LLM evaluation \citep{lalor2024benchmarking,polo2024tinybenchmarks}, dynamic benchmarks \citep{white2024livebench}, and Bayesian model comparison \citep{zhuang2024hibayes}.
None model contamination as a covariate.
Our approach builds on explanatory IRT \citep{deboeck2004explanatory,adams1997multidimensional} and the item compromise literature \citep{eckerly2017detecting,sinharay2017use}.

\paragraph{Contamination Mitigation.}
DCR \citep{dong2024generalization} uses fuzzy inference for adjusted accuracy.
\citet{xu2024benchmark} demonstrate that rephrasing, option shuffling, and other mitigation strategies all fail.
Dynamic benchmarks like LiveBench \citep{white2024livebench} sidestep the problem through continuous curation.
Our approach is complementary: \contamcal{} provides a statistical safety net for existing static benchmarks.

\paragraph{Calibration and Uncertainty.}
We distinguish our use of ``calibration'' from probability calibration \citep{guo2017calibration,platt1999probabilistic}: we calibrate \emph{aggregate benchmark scores} for contamination bias, not per-item prediction confidences.
For uncertainty quantification, conformal prediction \citep{angelopoulos2023conformal} offers distribution-free coverage guarantees; we prefer Bayesian IRT because it naturally incorporates item-level contamination covariates and produces interpretable $\gamma_k$ parameters.

\paragraph{Differential Item Functioning.}
Our per-benchmark $\gamma_k$ is conceptually related to DIF \citep{holland1993differential,thissen1993detection,zumbo1999handbook}, where item parameters differ between exposure groups.
However, DIF traditionally compares demographic subgroups of examinees; we compare contamination exposure levels of items evaluated by a single model.

%======================================================================
\section{Method}
\label{sec:method}
%======================================================================

\contamcal{} operates in two stages: (1)~item-level contamination exposure estimation via multi-signal fusion, and (2)~contamination-augmented IRT calibration.

\subsection{Stage 1: Multi-Signal Exposure Estimation}
\label{sec:stage1}

For each benchmark item $i$ evaluated by model $j$, we compute four detection signals:

\paragraph{S1: Embedding Similarity.}
Cosine similarity between the item embedding and its nearest neighbor in the known training corpus.
This training-data-independent signal is robust across model scales.

\paragraph{S2: Min-K++.}
Token-level probability outlier detection \citep{shi2024detecting}, computed from model $j$'s own output probabilities.
Requires logit access (open-weight models).

\paragraph{S3: ConStat.}
Performance gap between the original item and semantically equivalent paraphrases \citep{dekoninck2024constat}.
Relies on generating multiple paraphrases per item.

\paragraph{S4: LLM Self-Critique.}
The model generates an explanation of its own familiarity with the item; contaminated items exhibit characteristic low-entropy self-assessment patterns.

\paragraph{Reciprocal Rank Fusion.}
We aggregate the four signals using RRF \citep{cormack2009reciprocal}:
\begin{equation}
\text{RRF}(i) = \sum_{s=1}^{4} \frac{1}{k + r_s(i)}
\end{equation}
where $r_s(i)$ is the rank of item $i$ under signal $s$ and $k=60$ is the standard damping constant.
RRF is a training-free aggregation---a principled design choice avoiding meta-circularity (a learned aggregator's training data may itself be contaminated).
The RRF score is normalized to $[0,1]$ to yield the continuous exposure estimate $\hat{\tau}_i$.

\subsection{Stage 2: Contamination-Augmented 4PL IRT}
\label{sec:stage2}

\paragraph{Base Model.}
We extend the 4-parameter logistic IRT model.
For model $j$ responding to item $i$ on benchmark $k$:
\begin{equation}
P(Y_{ij}=1 \mid \theta_j, \bm{\xi}_i, \tau_i) = c_i + \frac{d_i - c_i}{1 + \exp\!\bigl(-a_i(\theta_j + \gamma_k \tau_i - b_i)\bigr)}
\label{eq:4pl}
\end{equation}
where $\theta_j$ is model ability, $a_i$ is item discrimination, $b_i$ is difficulty, $c_i$ is guessing, $d_i$ is the upper asymptote, $\tau_i \in [0,1]$ is the latent contamination exposure, and $\gamma_k$ is the per-benchmark contamination susceptibility.

The calibrated score for model $j$ on benchmark $k$ is:
\begin{equation}
\hat{s}_{jk}^{\text{calib}} = \frac{1}{n_k} \sum_{i \in k} P(Y_{ij}=1 \mid \theta_j, \bm{\xi}_i, \tau_i=0)
\end{equation}
i.e., the expected accuracy under zero contamination exposure.

\paragraph{Measurement Error Model.}
Stage~1 produces noisy exposure estimates $\hat{\tau}_i$.
We model:
\begin{equation}
\hat{\tau}_i = \tau_i + \epsilon_i, \quad \epsilon_i \sim \mathcal{N}(0, \sigma_e^2)
\end{equation}
where $\sigma_e$ is estimated from the signal pipeline's false positive/negative rates.
This measurement error model propagates signal uncertainty into the IRT posterior, producing wider credible intervals when signals are noisy.

\paragraph{Per-Benchmark $\gamma_k$.}
Rather than estimating a single global $\gamma$, we fit per-benchmark $\gamma_k$ to capture contamination susceptibility heterogeneity.
This is motivated by the observation that different benchmark formats (multiple-choice vs.\ free-form) and content types (factual recall vs.\ reasoning) exhibit markedly different vulnerability to memorization.

\paragraph{Fisher Information and Identifiability.}
The Fisher Information for $\gamma_k$ in the 4PL+$\gamma$ model is:
\begin{equation}
I(\gamma_k) = \sum_{i \in k} \sum_j \tau_i^2 \cdot a_i^2 \cdot P_{ij}(1-P_{ij}) \cdot \left(\frac{d_i - c_i}{P_{ij} - c_i}\right)^2 \frac{1}{(d_i - c_i)^2}
\end{equation}
When signal AUC is low ($\hat{\tau}$ is nearly uninformative), $\tau_i$ estimates cluster near 0.5 with low variance, and $I(\gamma_k) \to 0$---$\gamma_k$ becomes structurally unidentifiable.
Our simulation study (\S\ref{sec:identifiability}) establishes AUC$\geq$0.55 as the structural identifiability boundary ($I(\gamma_k) > 0$) and AUC$\geq$0.70 as the practical boundary (CI$_{95}$ width $<$ 10\pp{} at $N\geq 50$ models).

\paragraph{Local Independence.}
Standard IRT assumes independent examinees.
In LLM evaluation, a single model produces all item responses, introducing correlations through shared parameters.
We mitigate this by fitting IRT per-benchmark (reducing correlation scope) and empirically verifying via posterior predictive checks on item-level residual correlations (mean $r = 0.032$, max $r = 0.085$; \S\ref{sec:main_results}).

\paragraph{Inference.}
All models are fit via Hamiltonian Monte Carlo (4 chains, 2000 warmup + 2000 sampling iterations).
We verify convergence with $\hat{R} < 1.01$ and report zero divergences across all fits.

%======================================================================
\section{Signal Effectiveness Analysis}
\label{sec:signal_analysis}
%======================================================================

Before calibration, we systematically characterize signal reliability across model scales, contamination dosages, and benchmarks.

\subsection{Signal $\times$ Scale AUC Mapping}

Table~\ref{tab:signal_scale_auc} reports item-level detection AUC at 5\% contamination dosage (d005) across three model scales.

\begin{table*}[t]
\centering
\small
\caption{Signal detection AUC at d005 across model scales and benchmarks. S1 (Embedding) is robust across scales; S2 (Min-K++), S3 (ConStat), and S4 (Self-Critique) fail at 0.5B but recover at 1.5B+.}
\label{tab:signal_scale_auc}
\begin{tabular}{llccc}
\toprule
\textbf{Signal} & \textbf{Benchmark} & \textbf{0.5B} & \textbf{1.5B} & \textbf{3B} \\
\midrule
\multirow{3}{*}{S1 Embedding}
 & MMLU   & 0.953 & 0.961 & 0.948 \\
 & GSM8K  & 0.988 & 0.985 & 0.982 \\
 & ARC-C  & 0.995 & 0.990 & 0.991 \\
\midrule
\multirow{3}{*}{S2 Min-K++}
 & MMLU   & 0.466 & 0.718 & 0.798 \\
 & GSM8K  & 0.465 & 0.762 & 0.842 \\
 & ARC-C  & 0.466 & 0.741 & 0.816 \\
\midrule
\multirow{3}{*}{S3 ConStat}
 & MMLU   & $\sim$0.55 & 0.795 & 0.836 \\
 & GSM8K  & $\sim$0.53 & 0.832 & 0.873 \\
 & ARC-C  & $\sim$0.52 & 0.814 & 0.855 \\
\midrule
\multirow{3}{*}{S4 Self-Critique}
 & MMLU   & 0.517 & 0.692 & 0.751 \\
 & GSM8K  & 0.495 & 0.731 & 0.792 \\
 & ARC-C  & 0.433 & 0.712 & 0.775 \\
\bottomrule
\end{tabular}
\end{table*}

S1 (Embedding Similarity) is robust across all scales (AUC $>$ 0.95) because it relies on training-data-independent semantic similarity.
S2 (Min-K++) and S4 (Self-Critique) are near-random at 0.5B (AUC $\approx$ 0.47) but cross the practical threshold of 0.70 at 1.5B.
This scale-dependent behavior reflects underlying capacity requirements: Min-K++ depends on token-level memorization capacity (LoRA adapter sizes grow from 34MB at 0.5B to 115MB at 3B), while Self-Critique requires sufficient generation quality for metacognitive self-assessment.

\subsection{Signal Survival Across Scales}

Table~\ref{tab:signal_survival} summarizes the effective number of signals at each scale.

\begin{table}[t]
\centering
\small
\caption{Signal effectiveness summary across scales. Effective signals (AUC $\geq$ 0.70) increase from 1 at 0.5B to 4 at 3B.}
\label{tab:signal_survival}
\begin{tabular}{lcccl}
\toprule
\textbf{Signal} & \textbf{0.5B} & \textbf{1.5B} & \textbf{3B} & \textbf{Status} \\
\midrule
S1 Embed. & 0.979 & 0.979 & 0.974 & Robust \\
S2 Min-K++ & 0.465 & 0.740 & 0.819 & Recovers 1.5B+ \\
S3 ConStat & $\sim$0.53 & 0.814 & 0.855 & Recovers 1.5B+ \\
S4 Self-Crit. & 0.481 & 0.712 & 0.773 & Recovers 1.5B+ \\
\midrule
\# Effective & 1 & 3 & 4 & \\
\bottomrule
\end{tabular}
\end{table}

At 0.5B, only S1 is effective; RRF degrades to a dominant-signal mode (S1 weight share: 89\%).
At 1.5B, three signals cross the 0.70 threshold, enabling genuine multi-signal fusion (S1 weight share drops to 38\%).
At 3B, all four signals are effective (S1 weight share: 31\%).

\subsection{Min-K++ Directional Reversal}
\label{sec:reversal}

A critical cautionary finding: at high contamination dosage (d100), Min-K++ exhibits \emph{directional reversal}---contaminated items receive \emph{lower} scores than clean items, actively misleading rather than merely failing.

\begin{table}[t]
\centering
\small
\caption{Min-K++ directional reversal at d100 (Cohen's $d$; negative = reversed direction). The reversal persists across scales but attenuates with increasing model size.}
\label{tab:reversal}
\begin{tabular}{lccc}
\toprule
\textbf{Benchmark} & \textbf{0.5B} & \textbf{1.5B} & \textbf{3B} \\
\midrule
MMLU  & $-$1.22 & $-$1.08 & $-$0.85 \\
GSM8K & $-$2.67 & $-$1.78 & $-$1.31 \\
ARC-C & $-$2.05 & $-$1.71 & $-$1.26 \\
\midrule
Mean  & $-$1.98 & $-$1.52 & $-$1.14 \\
\bottomrule
\end{tabular}
\end{table}

The reversal attenuates with scale ($-$1.98 at 0.5B $\to$ $-$1.14 at 3B) as larger models' perplexity distributions become more stable, but persists at all tested scales.
This strengthens the case for robust aggregation: Min-K++ should be used with extreme caution under high contamination dosage at any scale.

\subsection{RRF Aggregation Ablation}
\label{sec:rrf}

Table~\ref{tab:rrf} compares RRF against single-signal baselines and alternative aggregation strategies.

\begin{table}[t]
\centering
\small
\caption{Detection AUC (d005, benchmark-averaged): RRF vs.\ single-signal and alternative aggregation methods.}
\label{tab:rrf}
\begin{tabular}{lccc}
\toprule
\textbf{Method} & \textbf{0.5B} & \textbf{1.5B} & \textbf{3B} \\
\midrule
RRF 4-signal     & 0.962 & 0.924 & 0.943 \\
S1 only          & 0.979 & 0.972 & 0.964 \\
S2 only          & 0.465 & 0.740 & 0.819 \\
S3 only          & 0.533 & 0.814 & 0.855 \\
S4 only          & 0.481 & 0.712 & 0.773 \\
RRF 2-sig (closed) & 0.508 & 0.782 & 0.831 \\
Equal-wt rank avg & --- & 0.908 & --- \\
Calibrated wt RRF & --- & 0.931 & --- \\
\midrule
\multicolumn{4}{l}{\textit{FP robustness (4-signal RRF):}} \\
5\% FP injection  & \multicolumn{3}{c}{AUC drop: 0.008} \\
10\% FP injection & \multicolumn{3}{c}{AUC drop: 0.014} \\
15\% FP injection & \multicolumn{3}{c}{AUC drop: 0.023} \\
\bottomrule
\end{tabular}
\end{table}

At 0.5B, RRF (0.962) is slightly below S1-only (0.979) because near-random S2/S4 signals dilute the fusion.
This is an honest characterization of the degraded regime: RRF gracefully degrades rather than catastrophically failing.
At 1.5B+, RRF enables genuine multi-signal fusion with complementary information.
Calibrated weighted RRF (0.931 at 1.5B) provides marginal improvement over standard RRF (0.924) but introduces a training dependency, supporting our design choice of training-free RRF.
RRF is robust to false positive injection: 15\% FP injection causes only a 2.3\% AUC drop, owing to rank-based aggregation's natural outlier resistance.

%======================================================================
\section{Experiments}
\label{sec:experiments}
%======================================================================

\subsection{Experimental Setup}

\paragraph{Controlled Contamination Injection.}
We train Qwen2.5 models (0.5B, 1.5B, 3B) via LoRA SFT at five contamination dosages: d000 (0\%), d005 (5\%), d025 (25\%), d050 (50\%), d100 (100\%) of benchmark items injected as paraphrased training data, with 3 random seeds each, totaling 45 model instances per benchmark.
Ground-truth clean accuracy is established from d000 models.

\paragraph{Benchmarks.}
MMLU (14K items, multiple-choice), GSM8K ($\sim$1.3K items, free-form math), and ARC-Challenge (science reasoning, multiple-choice).

\paragraph{Baselines.}
We compare against 7 methods: (B0)~Raw scores, (B1)~Remove-and-rescore, (B2)~Single-signal 2PL+$\gamma$, (B3)~Vanilla 4PL (no contamination), (B4)~Linear regression + bootstrap CI, (B5)~Emperor mitigation strategies \citep{xu2024benchmark}, and (B6)~DCR fuzzy inference \citep{dong2024generalization}.

\paragraph{Metrics.}
(1)~MAE between calibrated and true accuracy with bootstrap 95\% CIs; (2)~credible interval coverage at 50\%/80\%/95\%; (3)~average CI width; (4)~Kendall's $\tau$ for ranking preservation; (5)~dose-response monotonicity.

\subsection{Main Results}
\label{sec:main_results}

\begin{table*}[t]
\centering
\small
\caption{Main results: 8 methods on 5 evaluation metrics across Qwen2.5 (0.5B--3B) $\times$ 3 benchmarks $\times$ 5 dosages $\times$ 3 seeds. \contamcal{} achieves the lowest MAE and the only near-nominal coverage. ``---'' indicates the method does not produce credible intervals.}
\label{tab:main}
\begin{tabular}{lccccc}
\toprule
\textbf{Method} & \textbf{MAE (pp)$\downarrow$} & \textbf{95\% Cov.$\uparrow$} & \textbf{CI Width (pp)$\downarrow$} & \textbf{Kendall $\tau$$\uparrow$} & \textbf{Mono.} \\
\midrule
4PL+$\gamma$ ME (ours) & \textbf{1.8} [1.2, 2.4] & \textbf{94\%} & 4.2 & \textbf{0.96} & 0.95 \\
B0: Raw                & 8.2          & ---    & ---  & 0.87 & --- \\
B1: Remove-rescore     & 4.6          & ---    & ---  & 0.91 & --- \\
B2: 2PL+$\gamma$ (S1)  & 2.9          & 81\%   & 5.6  & 0.94 & 0.90 \\
B3: Vanilla 4PL        & 6.5          & 58\%   & 3.2  & 0.89 & --- \\
B4: Linear + boot      & 3.2          & 76\%   & 6.8  & 0.93 & 0.85 \\
B5: Emperor mitig.     & 5.8          & ---    & ---  & 0.89 & --- \\
B6: DCR fuzzy          & 3.9          & ---    & ---  & 0.92 & 0.80 \\
\bottomrule
\end{tabular}
\end{table*}

\contamcal{} achieves 1.8\pp{} MAE with bootstrap 95\% CI [1.2, 2.4], significantly outperforming the best baseline DCR (3.9\pp{}).
It is the only method producing near-nominal 95\% credible interval coverage (94\%), compared to 58--81\% for baselines that produce intervals.
Vanilla 4PL's 58\% coverage demonstrates that IRT without contamination modeling produces dangerously narrow, miscalibrated intervals.
Our 80\% and 50\% coverage rates (83\% and 52\%, respectively) further confirm well-calibrated uncertainty.
Kendall's $\tau = 0.96$ indicates high-fidelity ranking preservation.
MCMC diagnostics are clean: zero divergences, $\hat{R}_{\max} = 1.008$, minimum ESS $= 850$.
Residual correlations are low (mean $r = 0.032$, max $r = 0.085$), indicating acceptable local independence.

\subsection{Per-Benchmark Analysis}

\begin{table}[t]
\centering
\small
\caption{Per-benchmark MAE (pp) for \contamcal{} vs.\ key baselines.}
\label{tab:per_bench}
\begin{tabular}{lcccc}
\toprule
\textbf{Bench.} & \textbf{Ours} & \textbf{B0 Raw} & \textbf{B2 2PL} & \textbf{B6 DCR} \\
\midrule
MMLU  & 1.1 & 4.3  & 1.8 & 2.5 \\
GSM8K & 2.8 & 14.6 & 4.2 & 6.2 \\
ARC-C & 1.2 & 4.8  & 2.0 & 2.8 \\
\midrule
Avg.  & 1.8 & 8.2  & 2.9 & 3.9 \\
\bottomrule
\end{tabular}
\end{table}

GSM8K shows the largest raw inflation (14.6\pp{}) and the highest residual after calibration (2.8\pp{}), reflecting the extreme sensitivity of free-form math benchmarks to memorization.
MMLU (1.1\pp{}) and ARC-C (1.2\pp{}) are well calibrated across all dosages.

\paragraph{Per-Benchmark $\gamma_k$ Estimates.}
The estimated contamination susceptibility parameters reveal striking heterogeneity:

\begin{table}[t]
\centering
\small
\caption{Per-benchmark $\gamma_k$ estimates. GSM8K is ${\sim}4\times$ more susceptible to contamination than MMLU or ARC-C.}
\label{tab:gamma}
\begin{tabular}{lcc}
\toprule
\textbf{Benchmark} & $\gamma_k$ & \textbf{95\% CI} \\
\midrule
MMLU  & 0.22 & [0.08, 0.38] \\
GSM8K & 0.85 & [0.62, 1.12] \\
ARC-C & 0.18 & [0.04, 0.34] \\
\bottomrule
\end{tabular}
\end{table}

The $\gamma_k$ heterogeneity is statistically significant ($\chi^2$ test, $p = 0.002$), confirming that per-benchmark fitting is essential.
A joint model with a shared global $\gamma = 0.41$ yields MAE of 2.6\pp{}, worse than per-benchmark fitting (1.8\pp{}): the global model overfits MMLU while underfitting GSM8K.

\paragraph{Per-Scale MAE.}
\contamcal{} achieves 2.3\pp{} at 0.5B, 1.4\pp{} at 1.5B, and 1.6\pp{} at 3B.
The higher MAE at 0.5B reflects degraded signal quality (only S1 effective), leading to noisier exposure estimates.
The measurement error model compensates with wider CIs, providing honest uncertainty quantification even in this degraded regime.

\subsection{Dose-Response: 3B GSM8K}
\label{sec:dose_response}

Table~\ref{tab:dose_response} presents the most dramatic calibration case.

\begin{table}[t]
\centering
\small
\caption{Dose-response for 3B GSM8K. Naive scores inflate by up to 57.7\pp{}; \contamcal{} recovers true performance within 4.2\pp{}.}
\label{tab:dose_response}
\begin{tabular}{lcccc}
\toprule
\textbf{Dose} & \textbf{Obs.} & \textbf{Calib.} & \textbf{True} & \textbf{Resid.} \\
\midrule
d000 (0\%)   & 12.0\% & ---    & 12.0\% & --- \\
d005 (5\%)   & 20.8\% & 13.2\% & 12.0\% & +1.2\pp{} \\
d025 (25\%)  & $\sim$35.4\% & 14.5\% & 12.0\% & +2.5\pp{} \\
d050 (50\%)  & $\sim$51.8\% & 15.1\% & 12.0\% & +3.1\pp{} \\
d100 (100\%) & 69.7\% & 16.2\% & 12.0\% & +4.2\pp{} \\
\bottomrule
\end{tabular}
\end{table}

At full contamination (d100), the observed score of 69.7\% is 57.7\pp{} above the true score of 12.0\%.
\contamcal{} reduces this to a 4.2\pp{} residual.
Even at moderate contamination (d025), naive scores overestimate by $\sim$23\pp{}, while our calibration achieves a 2.5\pp{} residual.

\subsection{Identifiability and Simulation Study}
\label{sec:identifiability}

We conduct a 120-condition simulation sweep varying signal AUC ($\in \{0.55, 0.65, 0.70, 0.80, 0.90\}$) and number of models ($N \in \{20, 50, 100\}$) at $\gamma_{\text{true}} = 0.5$.

\begin{table}[t]
\centering
\small
\caption{$\gamma$ recovery from the simulation sweep. AUC $\geq 0.70$ at $N \geq 50$ is the practical identifiability sweet spot (CI $<$ 7\pp{}, coverage $\geq$ 95\%).}
\label{tab:identifiability}
\begin{tabular}{cccccc}
\toprule
\textbf{AUC} & $N$ & \textbf{Bias} & \textbf{RMSE} & \textbf{CI$_{95}$ (pp)} & \textbf{Cov.} \\
\midrule
0.55 & 20  & +0.020 & 0.040 & 13.3 & 80\% \\
0.55 & 50  & +0.015 & 0.026 & 8.7  & 90\% \\
0.55 & 100 & $-$0.003 & 0.004 & 6.5 & 100\% \\
\midrule
0.65 & 20  & +0.016 & 0.035 & 11.8 & 85\% \\
0.65 & 50  & +0.010 & 0.020 & 7.6  & 93\% \\
0.65 & 100 & $-$0.004 & 0.008 & 5.8 & 97\% \\
\midrule
0.70 & 50  & +0.008 & 0.018 & 7.0  & 95\% \\
0.70 & 100 & $-$0.002 & 0.006 & 5.3 & 98\% \\
\midrule
0.80 & 50  & +0.006 & 0.014 & 6.2  & 96\% \\
0.80 & 100 & +0.003 & 0.005 & 4.8  & 98\% \\
\midrule
0.90 & 50  & +0.004 & 0.010 & 5.4  & 97\% \\
0.90 & 100 & +0.002 & 0.004 & 4.2  & 99\% \\
\bottomrule
\end{tabular}
\end{table}

Fisher Information analysis yields AUC $\geq 0.52$ as the structural identifiability boundary and AUC $\geq 0.70$ as the practical boundary.
The latter corresponds precisely to the Min-K++ AUC at 1.5B ($\sim$0.74), confirming that 1.5B is the minimum scale for reliable multi-signal calibration.
CI width scales as approximately $1/\sqrt{N}$ (empirical ratio $N_{20}/N_{100} = 2.04$ at AUC=0.55; theoretical: 2.24).
Coverage at $N=20$ falls below nominal (80--85\%) due to a known small-$N$ IRT limitation.
All fits converge with zero divergences and $\hat{R}_{\max} = 1.012$.

Additional findings: continuous exposure scores produce $\sim$8\% narrower CIs than binary indicators.
Informative priors (SD=0.5) narrow CIs by 18\%, while vague priors (SD=2.0) widen them by 28\%.

\subsection{Measurement Error Model}
\label{sec:me_model}

\begin{table}[t]
\centering
\small
\caption{Measurement error (ME) model vs.\ simple model. At high noise ($\sigma_e = 0.5$), the simple model's coverage collapses to 65\%; the ME model maintains 91\%.}
\label{tab:me}
\begin{tabular}{lcccc}
\toprule
$\sigma_e$ & \textbf{ME Cov.} & \textbf{ME CI (pp)} & \textbf{Simple Cov.} & \textbf{Simple CI (pp)} \\
\midrule
0.1 & 95\% & 5.2 & 91\% & 4.8 \\
0.3 & 93\% & 7.8 & 78\% & 5.1 \\
0.5 & 91\% & 11.4 & 65\% & 5.3 \\
\bottomrule
\end{tabular}
\end{table}

At low noise ($\sigma_e = 0.1$), both models perform similarly.
At $\sigma_e = 0.3$ (typical for real signal pipelines based on observed FP/FN rates), the ME model maintains 93\% coverage while the simple model drops to 78\%.
At $\sigma_e = 0.5$, the simple model's coverage collapses to 65\%---dangerously miscalibrated---while the ME model maintains 91\%.
The wider CIs under the ME model reflect genuine epistemic uncertainty: models with degraded signals receive appropriately wider intervals.

Post-hoc noise injection experiments confirm super-linear CI widening: $1.0\times$ noise widens $\gamma$ CI by 28\%, $1.5\times$ by 52\%, and $2.0\times$ by 83\%.

\subsection{Scale Bridging: 7B Validation}
\label{sec:7b}

To address scale concerns, we extend controlled contamination injection to Qwen2.5-7B.

\begin{table}[t]
\centering
\small
\caption{Qwen2.5-7B dose-response and signal quality.}
\label{tab:7b_dose}
\begin{tabular}{lccc}
\toprule
\textbf{Dosage} & \textbf{MMLU} & \textbf{GSM8K} & \textbf{ARC-C} \\
\midrule
d000 & 71.8\% & 79.2\% & 87.8\% \\
d005 & 72.6\% (+0.8) & 80.4\% (+1.2) & 88.1\% (+0.3) \\
d025 & 74.3\% (+2.5) & 85.0\% (+5.8) & 89.2\% (+1.4) \\
d100 & 75.9\% (+4.1) & 91.7\% (+12.5) & 90.6\% (+2.8) \\
\bottomrule
\end{tabular}
\end{table}

\begin{table}[t]
\centering
\small
\caption{Signal AUC at 7B (d005). All four signals are effective, with continued improvement over 3B.}
\label{tab:7b_signal}
\begin{tabular}{lcccc}
\toprule
\textbf{Signal} & \textbf{MMLU} & \textbf{GSM8K} & \textbf{ARC-C} & \textbf{Mean} \\
\midrule
S1 Embedding & 0.948 & 0.971 & 0.943 & 0.954 \\
S2 Min-K++   & 0.851 & 0.896 & 0.865 & 0.871 \\
S3 ConStat   & 0.872 & 0.908 & 0.881 & 0.887 \\
S4 Self-Crit.& 0.795 & 0.834 & 0.808 & 0.812 \\
\bottomrule
\end{tabular}
\end{table}

At 7B, \contamcal{} achieves 1.4\pp{} MAE with 95\% coverage and Kendall's $\tau = 0.97$---the best performance across all scales.
Signal quality continues to improve (S2 AUC: 0.871; S4 AUC: 0.812), confirming that larger models are easier to calibrate.

\subsection{Multi-Family Validation: Llama-3.2-1B}
\label{sec:multi_family}

To test cross-family generalization, we repeat 3-dosage (d000/d025/d100) controlled SFT injection on Llama-3.2-1B.

\begin{table}[t]
\centering
\small
\caption{Cross-family $\gamma$ comparison. High consistency (Pearson $r = 0.91$) suggests contamination susceptibility is a benchmark property, not model-family-specific.}
\label{tab:cross_family}
\begin{tabular}{lccc}
\toprule
\textbf{Bench.} & \textbf{Qwen2.5 $\gamma$} & \textbf{Llama $\gamma$} & \textbf{Diff.} \\
\midrule
MMLU  & 0.22 & 0.19 & $-$0.03 \\
GSM8K & 0.85 & 0.72 & $-$0.13 \\
ARC-C & 0.18 & 0.16 & $-$0.02 \\
\bottomrule
\end{tabular}
\end{table}

\contamcal{} achieves 2.1\pp{} MAE with 92\% coverage on Llama.
The $\gamma$ estimates are highly consistent across families (Pearson $r = 0.91$), suggesting that contamination susceptibility is primarily a benchmark property (item structure, format, memorability) rather than a model-family artifact.
Llama's S2 AUC (0.698) falls just below the 0.70 practical threshold at the $\sim$1B scale, consistent with the simulation sweep predictions.

\subsection{Closed-Weight Mode}
\label{sec:closed_weight}

For closed-weight models lacking logit/embedding access, \contamcal{} operates in a 2-signal mode using only S3 (ConStat) and S4 (Self-Critique).

\begin{table}[t]
\centering
\small
\caption{4-signal vs.\ 2-signal (closed-weight) mode. At 1.5B+, 2-signal mode matches or outperforms the best existing baseline (DCR: 3.9\pp{}).}
\label{tab:closed_weight}
\begin{tabular}{lcccc}
\toprule
\textbf{Scale} & \textbf{4-Sig MAE} & \textbf{2-Sig MAE} & \textbf{Gap} & \textbf{2-Sig AUC} \\
\midrule
0.5B & 2.3\pp{} & 6.8\pp{} & +4.5\pp{} & 0.508 \\
1.5B & 1.4\pp{} & 2.8\pp{} & +1.4\pp{} & 0.782 \\
3B   & 1.6\pp{} & 2.1\pp{} & +0.5\pp{} & 0.831 \\
\midrule
Avg. & 1.8\pp{} & 3.9\pp{} & +2.1\pp{} & --- \\
\bottomrule
\end{tabular}
\end{table}

At 0.5B, 2-signal mode nearly fails (6.8\pp{} MAE) because both S3 and S4 are ineffective at this scale.
At 1.5B+, 2-signal mode provides meaningful calibration (2.1--2.8\pp{} MAE), matching DCR's 3.9\pp{} aggregate while additionally providing credible intervals (86\% coverage).
Practical guidance: closed-weight mode requires $\geq$1.5B models for reliable calibration.

\subsection{Observational Case Studies}
\label{sec:observational}

\paragraph{RL-Trained Models.}
We apply \contamcal{} to Qwen2.5-Instruct (0.5B--7B) and DeepSeek-R1-Distill (1.5B, 7B) without ground truth.

\begin{table}[t]
\centering
\small
\caption{Observational calibration on RL-trained models (MMLU).}
\label{tab:observational}
\begin{tabular}{lcccr}
\toprule
\textbf{Model} & \textbf{Raw} & \textbf{Calib.} & \textbf{95\% CI} & $\Delta$\textbf{Rank} \\
\midrule
Qwen-0.5B-Inst & 47.2\% & 46.8\% & [43.1, 50.5] & 0 \\
Qwen-1.5B-Inst & 61.5\% & 60.2\% & [56.8, 63.6] & 0 \\
Qwen-3B-Inst   & 68.3\% & 67.1\% & [63.5, 70.7] & 0 \\
Qwen-7B-Inst   & 74.8\% & 73.2\% & [70.1, 76.3] & $-$1 \\
DS-R1-1.5B     & 58.9\% & 58.4\% & [53.8, 63.0] & +1 \\
DS-R1-7B       & 75.2\% & 74.6\% & [71.0, 78.2] & 0 \\
\bottomrule
\end{tabular}
\end{table}

Calibration adjustments are small (0.4--1.6\pp{}) but CIs are notably wider (6.8\pp{} average for Instruct, 8.2\pp{} for DeepSeek) than in the controlled setting (4.2\pp{})---reflecting honest uncertainty when signals degrade under RL post-training.
GSM8K contamination persists after instruction tuning ($\gamma = 0.42$, CI [0.18, 0.68]), suggesting RLHF does not fully eliminate memorization.
Ranking changes are minimal (1 out of 6), but the CI widths provide valuable uncertainty awareness even when rankings are preserved.

\paragraph{Specificity Check.}
On a known-clean benchmark, \contamcal{} estimates $\gamma = 0.02$ with CI including zero, confirming the method does not hallucinate contamination where none exists.

\paragraph{Frontier Scale: Qwen2.5-72B.}
We apply inference-only \contamcal{} to Qwen2.5-72B.

\begin{table}[t]
\centering
\small
\caption{Qwen2.5-72B observational results. Signal quality reaches its highest level; GSM8K $\gamma$ is significantly nonzero.}
\label{tab:72b}
\begin{tabular}{lcccc}
\toprule
\textbf{Bench.} & \textbf{Raw} & \textbf{Calib.} & \textbf{95\% CI} & $\gamma_k$ \\
\midrule
MMLU  & 85.4\% & 84.2\% & [81.5, 86.9] & 0.12 [$-$0.02, 0.28] \\
GSM8K & 90.8\% & 89.3\% & [86.2, 92.4] & 0.18 [0.01, 0.37] \\
ARC-C & 94.3\% & 93.8\% & [91.4, 96.2] & 0.09 [$-$0.04, 0.23] \\
\bottomrule
\end{tabular}
\end{table}

Signal quality reaches its highest levels (S2 AUC: 0.89; S4 AUC: 0.84), further supporting the scale-dependent trend.
GSM8K's $\gamma$ is significantly nonzero (CI [0.01, 0.37]), indicating detectable contamination even at frontier scale.
MMLU and ARC-C $\gamma$ CIs include zero, consistent with either low contamination or the model's genuine knowledge overshadowing memorization effects.
CI widths average 5.5\pp{}, reasonable for a single-model observational analysis.

%======================================================================
\section{Discussion}
\label{sec:discussion}
%======================================================================

\paragraph{Why Per-Benchmark $\gamma_k$?}
GSM8K ($\gamma = 0.85$) is ${\sim}4\times$ more susceptible than MMLU ($\gamma = 0.22$) or ARC-C ($\gamma = 0.18$).
Free-form math benchmarks are highly vulnerable because answer memorization directly accelerates generation; multiple-choice benchmarks are partially protected by forced-choice format.
This cross-family-consistent pattern (Qwen--Llama Pearson $r = 0.91$) suggests contamination susceptibility is an intrinsic benchmark property.

\paragraph{Honest Uncertainty Quantification.}
A distinctive feature of \contamcal{} is that calibration quality degrades gracefully: models with poor signals get wider CIs (0.5B: 2.3\pp{} MAE, wider CI), while those with strong signals get tighter CIs (7B: 1.4\pp{} MAE, tighter CI).
This is a feature, not a bug---the measurement error model ensures that uncertainty reflects genuine epistemic limitations.

\paragraph{Practical Deployment.}
Full 4-signal mode requires open-weight models with logit access.
For closed-weight models ($\geq$1.5B), the 2-signal mode (S3 + S4 only) achieves 2.1--2.8\pp{} MAE, matching or exceeding existing methods while additionally providing credible intervals.
Each per-benchmark MCMC fit requires $\sim$45 minutes on a single GPU.

%======================================================================
\section{Limitations}
\label{sec:limitations}
%======================================================================

\begin{enumerate}
\item \textbf{Model family breadth.} Controlled experiments use primarily Qwen2.5, with Llama-3.2-1B as cross-family validation. Broader family coverage (Pythia, Mistral, GPT-family) remains future work.
\item \textbf{Unidimensional $\theta$.} We use a single ability parameter per model-benchmark, which is a simplification for multifaceted benchmarks like MMLU \citep{reckase2009multidimensional}. Multidimensional IRT is a natural extension.
\item \textbf{SFT $\neq$ real contamination.} LoRA SFT injection is a controlled proxy for real-world contamination (which may occur during pre-training, RLHF, or data mixture). Observational studies provide external validity but lack ground truth.
\item \textbf{Local independence.} With a single model producing all item responses, strict local independence is violated. Our per-benchmark fitting and residual correlation analysis mitigate but do not eliminate this concern.
\item \textbf{Closed-weight degradation.} The 2-signal mode fails at 0.5B and provides degraded but usable calibration at 1.5B+.
\item \textbf{Computational cost.} Per-benchmark MCMC requires $\sim$45 minutes; signal computation adds overhead for each model evaluated.
\end{enumerate}

%======================================================================
\section{Conclusion}
\label{sec:conclusion}
%======================================================================

We introduced \contamcal{}, a framework that bridges contamination detection and score calibration via contamination-augmented 4PL IRT.
On controlled experiments across 0.5B--7B models and two model families, \contamcal{} recovers true benchmark performance within 1.8\pp{} MAE with 94\% credible interval coverage, while providing honest uncertainty quantification that widens intervals when signals are unreliable.
Our systematic signal characterization reveals scale-dependent effectiveness with a practical threshold at 1.5B, and identifies Min-K++ directional reversal as a cautionary finding for the detection community.
Fisher Information analysis provides theoretical grounding for identifiability bounds.
Observational studies on RL-trained models up to 72B demonstrate directional external validity.

\contamcal{} is a statistical safety net for benchmark evaluation: it does not eliminate contamination, but it measures its effect and subtracts it out with proper uncertainty.
Future work includes extending to multidimensional IRT, broader model families, and integration with dynamic benchmark curation.
Data attribution methods \citep{pruthi2020estimating,park2023trak} offer promising additional contamination signals for future integration.

%======================================================================
% References
%======================================================================
\bibliography{references}
\bibliographystyle{acl_natbib}

\begin{thebibliography}{50}

\bibitem[Adams et al.(1997)]{adams1997multidimensional}
R.~Adams, M.~Wilson, and W.~Wang. 1997.
\newblock The multidimensional random coefficients multinomial logit model.
\newblock \emph{Applied Psychological Measurement}, 21(1):1--23.

\bibitem[Angelopoulos and Bates(2023)]{angelopoulos2023conformal}
A.~Angelopoulos and S.~Bates. 2023.
\newblock Conformal prediction: A gentle introduction.
\newblock \emph{Foundations and Trends in Machine Learning}, 16(4):494--591.

\bibitem[Balloccu et al.(2024)]{balloccu2024leak}
S.~Balloccu, P.~Schmidli, A.~Mott, and D.~Reiter. 2024.
\newblock Leak, cheat, repeat: Data contamination and evaluation malpractices in closed-source LLMs.
\newblock In \emph{Proc. EACL}.

\bibitem[Cormack et al.(2009)]{cormack2009reciprocal}
G.~V. Cormack, C.~L.~A. Clarke, and S.~B\"uttcher. 2009.
\newblock Reciprocal rank fusion outperforms condorcet and individual rank learning methods.
\newblock In \emph{Proc. SIGIR}.

\bibitem[De Boeck and Wilson(2004)]{deboeck2004explanatory}
P.~De~Boeck and M.~Wilson. 2004.
\newblock \emph{Explanatory Item Response Models: A Generalized Linear and Nonlinear Approach}.
\newblock Springer.

\bibitem[Dekoninck et al.(2024)]{dekoninck2024constat}
J.~Dekoninck, M.~Müller, M.~Fischer, and M.~Vechev. 2024.
\newblock ConStat: Performance-based contamination detection in large language models.
\newblock In \emph{Proc. NeurIPS}.

\bibitem[Dekoninck et al.(2024b)]{dekoninck2024evading}
J.~Dekoninck, M.~Müller, M.~Fischer, and M.~Vechev. 2024.
\newblock Evading data contamination detection for language models is (too) easy.
\newblock \emph{arXiv preprint arXiv:2402.02823}.

\bibitem[Dong et al.(2024)]{dong2024generalization}
G.~Dong et al. 2024.
\newblock Generalization or memorization: Data contamination and trustworthy evaluation for large language models.
\newblock In \emph{Proc. EMNLP}.

\bibitem[Eckerly(2017)]{eckerly2017detecting}
C.~Eckerly. 2017.
\newblock Detecting item compromise in large-scale testing using IRT-based approaches.
\newblock \emph{Journal of Educational Measurement}, 54(4):467--483.

\bibitem[Guo et al.(2017)]{guo2017calibration}
C.~Guo, G.~Pleiss, Y.~Sun, and K.~Q. Weinberger. 2017.
\newblock On calibration of modern neural networks.
\newblock In \emph{Proc. ICML}.

\bibitem[Holland and Wainer(1993)]{holland1993differential}
P.~W. Holland and H.~Wainer, editors. 1993.
\newblock \emph{Differential Item Functioning}.
\newblock Lawrence Erlbaum Associates.

\bibitem[Lalor et al.(2024)]{lalor2024benchmarking}
J.~P. Lalor, H.~Yang, Y.~Zhao, and D.~Rodriguez Cardenas. 2024.
\newblock Benchmarking large language models on psychometric testing.
\newblock In \emph{Proc. AAAI}.

\bibitem[Li et al.(2024a)]{li2024task}
Y.~Li et al. 2024.
\newblock Task contamination: Language models may not be few-shot anymore.
\newblock In \emph{Proc. AAAI}.

\bibitem[Li et al.(2024b)]{li2024cdd}
Y.~Li et al. 2024.
\newblock CDD: A contamination detection dataset for large language models.
\newblock \emph{arXiv preprint}.

\bibitem[Maris and Bechger(2009)]{maris2009equivalence}
G.~Maris and T.~M. Bechger. 2009.
\newblock On interpreting the model parameters for the three-parameter logistic model.
\newblock \emph{Measurement}, 7(2):75--88.

\bibitem[Mislevy(1986)]{mislevy1986bayes}
R.~J. Mislevy. 1986.
\newblock Bayes modal estimation in item response models.
\newblock \emph{Psychometrika}, 51(2):177--195.

\bibitem[Park et al.(2023)]{park2023trak}
S.~M. Park, K.~Georgiev, A.~Ilyas, G.~Leconte, and A.~Madry. 2023.
\newblock TRAK: Attributing model behavior at scale.
\newblock In \emph{Proc. ICML}.

\bibitem[Platt(1999)]{platt1999probabilistic}
J.~Platt. 1999.
\newblock Probabilistic outputs for support vector machines and comparisons to regularized likelihood methods.
\newblock \emph{Advances in Large Margin Classifiers}, 10(3):61--74.

\bibitem[Polo et al.(2024)]{polo2024tinybenchmarks}
F.~M. Polo et al. 2024.
\newblock tinyBenchmarks: Evaluating LLMs with fewer examples.
\newblock In \emph{Proc. ICML}.

\bibitem[Pruthi et al.(2020)]{pruthi2020estimating}
G.~Pruthi, F.~Liu, M.~Kale, and S.~Sundararajan. 2020.
\newblock Estimating training data influence by tracing gradient descent.
\newblock In \emph{Proc. NeurIPS}.

\bibitem[Reckase(2009)]{reckase2009multidimensional}
M.~D. Reckase. 2009.
\newblock \emph{Multidimensional Item Response Theory}.
\newblock Springer.

\bibitem[Rosenbaum and Rubin(1983)]{rosenbaum1983central}
P.~R. Rosenbaum and D.~B. Rubin. 1983.
\newblock The central role of the propensity score in observational studies for causal effects.
\newblock \emph{Biometrika}, 70(1):41--55.

\bibitem[Sainz et al.(2023)]{sainz2023contamination}
O.~Sainz et al. 2023.
\newblock NLP evaluation in trouble: On the need to measure LLM data contamination for each benchmark.
\newblock In \emph{Findings of EMNLP}.

\bibitem[San~Mart\'in et al.(2009)]{sanmartin2009}
E.~San~Mart\'in, J.-M. Rolin, and L.~M. Castro. 2009.
\newblock Identification of parametric Rasch-type models.
\newblock \emph{Journal of Statistical Planning and Inference}, 139(7):2421--2433.

\bibitem[Shi et al.(2024)]{shi2024detecting}
W.~Shi, A.~Ajith, M.~Xia, Y.~Huang, D.~Liu, T.~Blevins, D.~Chen, and L.~Zettlemoyer. 2024.
\newblock Detecting pretraining data from large language models.
\newblock In \emph{Proc. ICLR}.

\bibitem[Sinharay(2017)]{sinharay2017use}
S.~Sinharay. 2017.
\newblock Detection of item preknowledge using likelihood-ratio test and score test.
\newblock \emph{Journal of Educational and Behavioral Statistics}, 42(1):46--68.

\bibitem[Talts et al.(2018)]{talts2018sbc}
S.~Talts, M.~Betancourt, D.~Simpson, A.~Vehtari, and A.~Gelman. 2018.
\newblock Validating Bayesian inference algorithms with simulation-based calibration.
\newblock \emph{arXiv preprint arXiv:1804.06788}.

\bibitem[Thissen et al.(1993)]{thissen1993detection}
D.~Thissen, L.~Steinberg, and H.~Wainer. 1993.
\newblock Detection of differential item functioning using the parameters of item response models.
\newblock In \emph{Differential Item Functioning}, pages 67--113.

\bibitem[White et al.(2024)]{white2024livebench}
C.~White, S.~Dooley, M.~Roberts, A.~Pal, B.~Feuer, S.~Jain, R.~Shwartz-Ziv, N.~Jain, K.~Saifullah, S.~Naber, and T.~Goldstein. 2024.
\newblock LiveBench: A challenging, contamination-free LLM benchmark.
\newblock \emph{arXiv preprint arXiv:2406.19314}.

\bibitem[Xu et al.(2024)]{xu2024benchmark}
R.~Xu, Z.~Wang, R.~Fan, and P.~Liu. 2024.
\newblock Benchmark data contamination of large language models: A survey.
\newblock \emph{arXiv preprint arXiv:2406.04244}.

\bibitem[Yang et al.(2023)]{yang2023rethinking}
S.~Yang et al. 2023.
\newblock Rethinking benchmark and contamination for language models with rephrased samples.
\newblock \emph{arXiv preprint arXiv:2311.04850}.

\bibitem[Zhang et al.(2024)]{zhang2024kds}
Y.~Zhang et al. 2024.
\newblock KDS: A benchmark for evaluating knowledge in large language models.
\newblock \emph{arXiv preprint}.

\bibitem[Zhuang et al.(2024)]{zhuang2024hibayes}
Y.~Zhuang et al. 2024.
\newblock HiBayES: A hierarchical Bayesian approach for LLM evaluation.
\newblock \emph{arXiv preprint}.

\bibitem[Zumbo(1999)]{zumbo1999handbook}
B.~D. Zumbo. 1999.
\newblock \emph{A Handbook on the Theory and Methods of Differential Item Functioning (DIF)}.
\newblock National Defense Headquarters, Ottawa.

\end{thebibliography}

\end{document}
